% §5 Harness Architecture — target 2 pages
% Purpose: describe the method. Emphasize: network-native, parallelization, tractable context.
% Key figure: Fig 2 — pipeline 5 phases; Fig 3 — per-device parallelization.

\section{\harnessname: A Network-Native LLM Pentest Harness}
\label{sec:harness}

% ── §5.1 overview ─────────────────────────────────────────────────────────
\subsection{Overview}
\label{sec:harness:overview}

\begin{figure*}[t]
  \centering
  % \includegraphics[width=0.9\linewidth]{figures/fig2_pipeline.pdf}
  \fbox{\parbox{0.9\linewidth}{\centering\textit{
    Fig.~2 placeholder: 5-phase pipeline diagram.\\
    Phase 1 (topology prior) $\to$ Phase 2 (discovery) $\to$ Phase 3 (per-device triage)
    $\to$ Phase 4 (per-vuln verification) $\to$ Phase 5 (network report).
  }}}
  \caption{\harnessname{} pipeline. Solid arrows are deliverables; dashed boxes
  indicate LLM agent calls; Phase~3 and Phase~4 spawn micro-agents in parallel
  (§\ref{sec:harness:phase3}--§\ref{sec:harness:phase4}).}
  \label{fig:pipeline}
\end{figure*}

The harness decomposes network pentest into five phases with typed deliverables.
The key design decision is \emph{parallelization at the device and
vulnerability granularity}: rather than a single monolithic agent that must hold
the full network state, we spawn dedicated micro-agents per target with
constrained tool sets, then aggregate deterministically. This keeps LLM context
bounded at $O(1)$ per device regardless of network size.

% ── §5.2 Phase 1 ──────────────────────────────────────────────────────────
\subsection{Phase~1: Topology prior}
\label{sec:harness:phase1}

Phase~1 loads a graph model of the target network as a directed graph
$G=(V,E)$ where $V$ are devices and $E$ are links labeled by protocol. Two
modes are supported: (i)~\emph{scenario mode} loads the predefined topology
from the benchmark YAML; (ii)~\emph{discovery mode} starts from an empty graph
and instructs Phase~2 to populate $V$ via network scans. The output is a
markdown deliverable enumerating attack surface and preliminary pivot candidates
(betweenness-centrality ranked).

% ── §5.3 Phase 2 ──────────────────────────────────────────────────────────
\subsection{Phase~2: Full-network discovery}
\label{sec:harness:phase2}

Phase~2 executes active reconnaissance across the target \texttt{/24}:
\texttt{nmap\_discovery} (host discovery), \texttt{nmap} (service/version),
\texttt{arp\_scan} (L2 + vendor), and protocol-specific probes
(\texttt{mqtt\_listen} subscribing to \texttt{\#} on discovered brokers,
\texttt{ssh\_audit} against SSH servers, \texttt{curl\_headers} on HTTP).
Unlike single-host harnesses, the agent is not given a target IP — it must
derive one from the CIDR scope.

% ── §5.4 Phase 3 ──────────────────────────────────────────────────────────
\subsection{Phase~3: Per-device parallel triage}
\label{sec:harness:phase3}

\begin{figure}[t]
  \centering
  % \includegraphics[width=\linewidth]{figures/fig3_parallelization.pdf}
  \fbox{\parbox{0.95\linewidth}{\centering\textit{
    Fig.~3 placeholder: Phase 3 parallelization — $n$ devices $\to$ $n$ LLM agents
    in parallel $\to$ deterministic merge.
  }}}
  \caption{Per-device parallelization. Each micro-agent receives only its
  device's scan output, avoiding the context explosion of a monolithic
  network-level agent.}
  \label{fig:parallel}
\end{figure}

Phase~3 runs in three sub-phases:
\begin{enumerate}
  \item \textbf{Deterministic scanner pass.} Ansible-driven subprocess
  wrappers (\texttt{ssh\_audit}, \texttt{curl\_headers}, \texttt{mqtt\_listen},
  \texttt{modbus\_scan}, \ldots) run in parallel per discovered device and
  dump structured scan results to disk.
  \item \textbf{Per-device LLM triage.} For each device, a dedicated LLM agent
  receives only that device's Phase~2+3a outputs and has access to a minimal
  tool set: \texttt{http\_get} (follow-up file retrieval) and
  \texttt{cve\_search} (CPE-to-CVE lookup via NVD). The agent produces a JSON
  vuln list scoped to that device.
  \item \textbf{Deterministic merge.} Per-device outputs are concatenated,
  canonicalized through the taxonomy (\S\ref{sec:harness:taxonomy}), and
  deduplicated via severity-aware rules.
\end{enumerate}

This decomposition is essential for network scale: a monolithic agent on a
15-device scenario would need to fit $\sim$30--80k tokens of scan output into
a single context and reason about cross-device relationships
simultaneously. Per-device micro-agents flatten this to a constant size, at
the cost of losing some cross-device inference capability, which is recovered
in Phase~5.

% ── §5.5 Phase 4 ──────────────────────────────────────────────────────────
\subsection{Phase~4: Per-vulnerability verification with evidence levels}
\label{sec:harness:phase4}

Phase~4 spawns one micro-agent per Phase~3 finding (in parallel, with a
concurrency cap). Each micro-agent receives the finding (IP, type, port,
details) and full operational tool access: \texttt{ssh\_login},
\texttt{mysql\_query}, \texttt{mqtt\_listen}, \texttt{http\_get},
\texttt{ftp\_list}, \texttt{modbus\_scan}, \texttt{telnet\_connect}. The agent
is prompted to produce the deepest evidence level it can (L1--L3,
\S\ref{sec:iotbench:eval}) within a bounded turn budget.
Outputs include: \texttt{status} (\textsc{confirmed}, \textsc{not\_exploitable},
\textsc{error}), \texttt{evidence\_level}, \texttt{command} executed,
\texttt{stdout} excerpts, and \texttt{data\_extracted} (when L3).

Findings in \textsc{confirmed} state from Phase~3 override Phase~4
\textsc{error}/\textsc{failed} statuses (rationale:
Phase~3 is hypothesis generation with directly-observed indicators; Phase~4
crashes should not erase directly-observed findings).

% ── §5.6 Phase 5 ──────────────────────────────────────────────────────────
\subsection{Phase~5: Network-level report}
\label{sec:harness:report}

Phase~5 consumes all prior deliverables and produces a network-scoped report
with: device inventory, per-device vulnerabilities with severity, attack
surface summary, and actionable recommendations ordered by severity $\times$
exposure. This is the only phase with cross-device reasoning, deliberately
placed at the end to avoid bloating earlier phases.

% ── §5.7 taxonomy & tool abstraction ──────────────────────────────────────
\subsection{Taxonomy and tool abstraction}
\label{sec:harness:taxonomy}

A single module resolves LLM-produced vulnerability type strings into a
canonical taxonomy of 11~categories (\cref{tab:categories}), filtering noise
(e.g., meta-observations) and CONFIG-only findings. Tools are defined
declaratively as YAML (schema: name, description, JSON-schema parameters,
subprocess command template), so adding a new protocol probe requires no
Python code.

% ── §5.8 complexity ───────────────────────────────────────────────────────
\subsection{Complexity and cost envelope}
\label{sec:harness:cost}

For a network with $n$ devices and an average of $v$ findings per device, the
harness makes $O(1+n+nv)$ LLM calls (Phases 1, 3b, 4 respectively), bounded
per call by fixed turn budgets ($50/50/10/10/20$ for phases 1--5). Total cost
scales linearly with network size, not combinatorially.
